% 本文件由 LaTeX 排版 Agent 根据有效 translation.json 自动生成。
% author=Liturgies; work_id=0717; title=圣雅各伯神圣礼仪

\chapter{\WorkTitle{圣雅各伯神圣礼仪}}

% structure: p0001 source=h1 target=omitted reason=page-title-already-rendered-by-parent

% structure: p0002 source=h2 target=section reason=relative-heading-rank; base=h2

\section{1}

\Emph{司铎。}\par

我，啊，\Emph{主权者}主，我们的天主，求祢不要弃绝我，因我沾染了许多罪过；因为，请看，我来参与祢这神圣而天上的奥迹，并非因为堪当，而是唯靠祢的良善，我向祢发出呼声：天主，可怜我这个罪人；我得罪了天，也得罪了祢，不配来到祢这神圣而属灵祭台前，在这祭台上，祢的独生子、我们的主耶稣基督，为我这罪人、满身污秽者，奥迹地作为祭献而陈列。因此，我向祢献上这祈求与感恩，求祢派遣祢的圣神——安慰者——降临于我，坚强并装备我以胜任这服务；并使我得以不受谴责地向子民宣明由祢藉我所传的圣言，因我们的主耶稣基督，祢与祂及祢的全圣、良善、赋予生命、同体的圣神同受赞颂，自今至永远，直至无穷之世。阿们。\par

\Emph{立于祭台旁的祈祷。}\par

II 光荣归于圣父、圣子及\OriginalTerm{天主圣三}{Trinitas}的\OriginalTerm{三一神性之光}{triune light}，这光是在\OriginalTerm{三位}{trinitate}中合而为一，分开而不可分割的：因为天主圣三是唯一全能的天主，诸天宣扬祂的光荣，大地传述祂的统治，海洋彰显祂的大能，一切有知觉和理性的受造物时时颂扬祂的威严：因为一切光荣、尊威、权能、伟大和威严都属于祂，自今至永远，直至无穷之世。阿们。\par

\Emph{开始时的献香祈祷。}\par

III \OriginalTerm{主权者}{Sovereign}主耶稣基督，\OriginalTerm{天主的圣言}{\GreekText{Λόγος} \GreekText{τοῦ} \GreekText{Θεοῦ}}，祢曾甘愿在十字架上向天主圣父献上自己为\OriginalTerm{无玷的祭献}{\GreekText{ἄμωμος} \GreekText{θυσία}}，那具有\OriginalTerm{二性}{\GreekText{δύο} \GreekText{φύσεις}}的火炭，曾用来钳触及先知的口唇，除掉他的罪过，求祢也触及我们罪人的心，洁净我们的一切污秽，使我们圣洁地站在祢的圣祭台前，好向祢献上赞颂的祭献；并求祢接纳我们这些无用的仆人，将这馨香的乳香献上，熏香我们灵魂与身体的恶臭，并以祢全圣之神的圣化大能洁净我们：因为唯有祢是圣的，祢圣化人，并与信众相通；光荣归于祢，与祢的永恒圣父，及祢的全圣、良善、赋予生命的圣神，自今至永远，直至无穷之世。阿们。\par

\Emph{开始的祈祷。}\par

IV 仁善的永恒君王、宇宙的创造者，求祢接纳祢的教会，她藉着祢的基督来到祢面前；求祢赐予每人有益之事，引导众人达于成全，并使我们完全堪当祢圣化的恩宠，聚集我们在祢的圣教会内，这教会是祢用祢独生子、我们的主和救主耶稣基督的宝血所赎买的，祢与祂及祢的全圣、良善、赋予生命的圣神同受赞颂和荣耀，自今至永远，直至无穷之世。阿们。\par

\Emph{执事。}\par

V 让我们再次祈求主。\par

\Emph{司铎，会众进堂时的献香祈祷。}\par

天主，祢曾接纳亚伯尔的献礼、诺厄和亚巴郎的祭献、亚郎和匝加利亚的乳香，求祢也接纳我们罪人手中所献的这乳香，当作馨香的香气，并为赦免我们的罪过及祢全体子民的罪过；因为祢是应受赞颂的，光荣属于祢，父、子及圣神，自今至永远。\par

\Emph{执事。}\par

主上，请赐祝福。\par

\Emph{司铎祈祷。}\par

我们的主和天主耶稣基督，祢因无比的良善和不可抑制的爱，曾被钉在十字架上，又甘愿被长矛和钉子刺透；祢为我们设立了这\OriginalTerm{神秘可畏的礼仪}{\GreekText{μυστήριον} \GreekText{φρικτόν}}作为永久的纪念：求祢祝福祢在基督天主内的服务，祝福我们的进堂，并以祢不可言喻的慈悲，圆满完成我们这服务的呈献，自今至永远，直至无穷之世。阿们。\par

\Emph{执事的应祈祷。}\par

VI 愿主祝福我们，并使我们堪当如同色辣芬般呈献礼物，并咏唱神圣\OriginalTerm{三圣颂}{\GreekText{τρισάγιον}}的不断颂歌，藉着一切圣德成全的丰满和极其丰盛，自今至永远。\par

\Emph{然后执事开始咏唱进堂咏。}\par

\OriginalTerm{独生子}{\GreekText{μονογενής} \GreekText{Υἱός}}和\OriginalTerm{天主的圣言}{\GreekText{Λόγος} \GreekText{τοῦ} \GreekText{Θεοῦ}}，祢是\OriginalTerm{不朽的}{\GreekText{ἀθάνατος}}；祢为我们的救恩，甘愿取了圣\OriginalTerm{天主之母}{\GreekText{Θεοτόκος}}、\OriginalTerm{终身童贞}{\GreekText{ἀειπάρθενος}}玛利亚的肉身，不变地成了人并被钉在十字架上，基督我们的天主，祢以死亡踏碎了死亡；祢是\OriginalTerm{天主圣三}{Trinitas}中的一位，与父及圣神同受荣耀，求祢拯救我们。\par

\Emph{司铎从大门到祭台前诵念此祈祷。}\par

VII 全能的天主，大能荣耀的主，祢藉着祢独生子、我们的主、天主和救主耶稣基督在人间居住，赐给了我们进入\OriginalTerm{至圣所}{\GreekText{Ἅγια} \GreekText{Ἁγίων}}的入口，我们恳求和呼求祢的良善，因为我们即将站立在祢的圣祭台前，心怀恐惧和战栗；求祢，天主，向我们派遣祢的仁慈恩宠，圣化我们的灵魂、身体和心神，引领我们的思念归向虔敬，好使我们在纯洁的良心中，向祢呈献礼物、祭品和果实，以赦免我们的过犯，并为祢全体子民作\OriginalTerm{赎罪祭}{\GreekText{ἱλασμός}}，藉着祢独生子的恩宠、慈悲和仁爱，祢与祂同受赞颂，直至永远。阿们。\par

\Emph{走近祭台后，司铎说：——}\par

VIII 愿平安与众人。\par

\Emph{群众。}\par

也与你的心灵。\par

\Emph{司铎。}\par

愿主祝福我们众人，并圣化我们，使我们能参与并举行神圣而纯洁的\OriginalTerm{奥迹}{\GreekText{μυστήρια}}，使蒙福的灵魂在善人和义人之中获得安息，藉着祂的恩宠和仁爱，自今至永远，直至无穷之世。阿们。\par

\Emph{然后执事诵念祷文。}\par

IX 在平安中让我们祈求主。\par

为从上而来的平安、为天主对人的爱，并为我们灵魂的救恩，让我们祈求主。\par

为全世界的和平、为所有神圣教会的合一，让我们祈求主。\par

为赦免我们的罪、宽恕我们的过犯，并为我们能脱离一切困苦、忿怒、危险、艰难及仇敌的攻击，让我们祈求主。\par

\Emph{然后歌咏者咏唱三圣颂。}\par

圣哉天主，圣哉大能者，圣哉不朽者，求祢垂怜我们。\par

\Emph{然后司铎俯首祈祷。}\par

十、仁慈而怜悯、延缓发怒、极富恩宠而真实的天主，求祢从祢预备的居所垂顾，俯听我们这些恳求者，救我们脱离魔鬼和人的一切诱惑；不要扣留祢的援助，也不要以超过我们力量的惩罚加于我们：因为我们不能胜过那敌对者；但祢，\Person{主}，能救我们脱离一切反对我们的事。天主，求祢按祢的良善，救我们脱离这世界的困难，好使我们以纯洁的良心就近祢的圣祭台，我们就能与天上的万军一同无咎地向祢献上蒙福的\OriginalTerm{三圣颂}{Trisagion}赞歌，并因完成了悦乐祢的神圣礼仪，得以相称于永生。\par

\Emph{(高声)}\par

因为祢是圣的，\Person{主}我们的天主，祢居住在圣所中，我们向祢，圣父、圣子及圣神，献上赞美和\OriginalTerm{三圣颂}{Trisagion}赞歌，从今直到永远。\par

\Emph{会众。}\par

阿们。\par

\Emph{司铎。}\par

十一、愿平安与众人同在。\par

\Emph{会众。}\par

也与你的心灵同在。\par

\Emph{歌咏团。}\par

阿肋路亚。\par

随后按顺序宣读旧约和先知的神圣圣言；并宣告天主圣子的降生成人、祂的苦难、从死者中复活、升天，以及祂荣耀的再度来临；此事每日在神圣的礼仪中举行。\par

\Emph{宣读和训导后，执事说：——}\par

十二、我们大家说：主，求你垂怜。\par

全能的主，我们祖先的天主；\par

我们恳求祢，俯听我们。\par

为上天的和平，并为拯救我们的灵魂；\par

让我们恳求主。\par

为全世界的和平，以及天主所有圣教会的合一；\par

让我们恳求主。\par

为所有爱基督之人的得救和帮助；\par

我们恳求祢，俯听我们。\par

为我们能脱离一切苦难、忿怒、危险、困厄、被掳、惨死和我们的罪愆；\par

我们恳求祢，俯听我们。\par

为环绕站立、等候祢丰盛而充溢仁慈的百姓；\par

我们恳求祢，求祢仁慈和施恩。\par

主，求祢拯救祢的百姓，祝福祢的基业。\par

求祢以仁慈和怜悯眷顾祢的世界。\par

赖宝贵而赋予生命的十字架的大能，高举基督徒的角。\par

最仁慈的主，我们恳求祢，俯听我们向祢祈祷，并求祢垂怜我们。\par

\Emph{会众（三次）。}\par

主，求祢垂怜我们。\par

\Emph{执事。}\par

十三、为赦免我们的罪，宽恕我们的过犯，并使我们脱离一切苦难、忿怒、危险和困厄，让我们恳求主。\par

让我们向主恳求，使我们能度过完美、圣洁、平安而无罪的一整天。\par

让我们向主恳求一位和平的使者、忠实的向导、我们灵魂和身体的守护者。\par

让我们向主恳求宽恕和赦免我们的罪愆和过犯。\par

让我们向主恳求对我们灵魂有益和适当的事物，并为世界祈求和平。\par

让我们向主恳求，使我们能在平安与健康中度完余下的时光。\par

让我们恳求，使我们生命的终结是基督徒式的，没有痛苦、没有羞耻，并在基督可畏而威严的审判台前有好的辩护。\par

\Emph{司铎。}\par

十四、因为祢是福音与光明，我们灵魂和身体的救主与守护者，天主，以及祢的独生子，和祢至圣的圣神，从今直到永远。\par

\Emph{会众。}\par

阿们。\par

\Emph{司铎}\par

天主，祢教导了我们祢神圣而拯救的圣言，求祢光照我们罪人的灵魂，使我们能理解之前所宣讲的事，好让我们不仅被视为属灵事物的听者，也成为善行的实行者，追求无伪的信德、无瑕的生活和纯洁的交往。\par

\Emph{(高声)}\par

在基督耶稣我们的主内，祂与祢同受赞颂，连同祢至圣、良善而赋予生命的圣神，从今直到永远。\par

\Emph{会众。}\par

阿们。\par

\Emph{司铎。}\par

十五、愿平安与众人同在。\par

\Emph{会众。}\par

也与你的心灵同在。\par

\Emph{执事。}\par

让我们向主俯首。\par

\Emph{会众。}\par

向祢，主。\par

\Emph{司铎祈祷说：}\par

主宰的生命赐予者、美善的供应者，祢赐给了人类永生蒙福的希望，我们的主耶稣基督，求祢使我们圣洁相称，并完善这神圣礼仪，以享将来的真福。\par

\Emph{(高声)}\par

好使我们时常蒙祢的能力护佑，并被引领进入真理之光，就能向祢，圣父、圣子及圣神，献上赞美和感恩，从今直到永远。\par

\Emph{会众。}\par

阿们。\par

\Emph{执事。}\par

十六、不要让慕道者留下，不要让未受洗者留下，不要让不能与我们一同祈祷的人留下。你们彼此相看。那门。\par

全体端正：让我们再次向主祈祷。\par

% structure: p0106 source=h2 target=section reason=relative-heading-rank; base=h2

\section{2}

\Emph{司铎诵念香的祷文。}\par

全能的主宰，荣耀的君王，祢在万物受造以前即知晓一切，求祢在这神圣的时刻向我们这些呼求祢的人显现，救我们脱离过犯的羞耻；洁净我们的心思和意念，除去不洁的欲望、世俗的欺骗和魔鬼的一切影响；并应从我们罪人手中接受这香，如同祢接受了\Person{亚伯尔}、\Person{诺厄}、\Person{亚郎}、\Person{撒慕尔}及祢所有圣人的奉献一样，保护我们远离一切邪恶，使我们不断悦乐、朝拜和荣耀祢，圣父，及祢的独生子，和祢至圣的圣神，从今直到永远。\par

\Emph{读经员开始唱革鲁宾赞歌。}\par

愿所有血肉静默，以敬畏和战栗站立，心中不要思想任何尘世之事：——\par

因为万王之王、万主之主、我们的天主\Person{基督}前来被祭献，并作为食物赐给信众；天使的队列在祂前面行进，包括一切权能和宰制，多目的\OriginalTerm{革鲁宾}{Cherubim}和六翼的\OriginalTerm{色辣芬}{Seraphim}，遮住面容，高声呼喊：阿肋路亚，阿肋路亚，阿肋路亚。\par

\Emph{司铎携圣礼进来时，诵念此祷文：——}\par

XVII. 天主，我们的天主，祢派遣了天上的食粮，全世界的养料，我们的主耶稣基督，作为救主、救赎者、恩主，祝福我们并圣化我们，求祢亲自祝福这奉献，并仁慈地收纳它到祢天上的祭台：\par

求祢以良善和爱心纪念那些奉献的人和那些为他们奉献的人，并使我们无可指责地事奉祢神圣的奥迹：因为祢至可敬、伟大的圣名——圣父、圣子、圣神——受颂扬和光荣，自今至永远，及世之世。\par

\Emph{司铎：}\par

愿平安赐与众位。\par

\Emph{执事：}\par

主，请颁赐祝福。\par

\Emph{司铎：}\par

愿受颂扬的是天主，祂在呈献神圣而纯洁的奥迹时祝福并圣化我们众人，并使有福的灵魂安息于圣者义人之中，自今至永远，及世之世。\par

\Emph{执事：}\par

XVIII 让我们明智地留意。\par

\Emph{司铎开始：}\par

我信唯一的天主，全能的圣父，天地万物的创造者；并信唯一的主耶稣基督，天主子；以及信经的其余部分。\par

\Emph{然后他俯首祈祷：}\par

XIX. 万有的天主和主宰，求祢使我们这些不配的人堪当此一时辰，爱人者；使我们脱离一切诡诈和伪善，藉着和平与爱德的纽带彼此结合，因祢独生子、我们的主和救主耶稣基督所赐的祢神性知识的圣化而得以坚固；祢与祂同受颂扬，偕同祢至圣、良善、赋予生命的圣神，自今至永远，及世之世。阿们。\par

\Emph{执事：}\par

XX. 让我们端正地站立，虔诚地站立，怀着对天主的敬畏和心灵的痛悔站立。让我们在平安中向主祈祷。\par

\Emph{司铎：}\par

因为祢是和平、慈悲、仁爱、怜悯和慈祥的天主，祢的独生子，以及祢至圣的圣神，自今至永远。\par

\Emph{众：}\par

阿们。\par

\Emph{司铎：}\par

愿平安赐与众位。\par

\Emph{众：}\par

也与你的心灵同在。\par

\Emph{执事：}\par

让我们以圣洁的亲吻彼此致意。让我们向主俯首。\par

\Emph{司铎俯首，念以下祷文：}\par

XXI. 唯一的主和仁慈的天主，求祢向那些在祢圣洁的祭台前俯首，并寻求从祢而来的神性恩赐的人，赐下祢的善意恩宠；并以一切不可剥夺的神性祝福祝福我们众人，祢居住在高天，却垂顾卑微之事。\par

\Emph{（高声）}\par

因为祢至圣的圣名——圣父、圣子、圣神——堪受赞美、钦崇和至高的光荣，自今至永远，及世之世。\par

\Emph{执事：}\par

主，请颁赐祝福。\par

\Emph{司铎：}\par

主要祝福我们，并藉祂的恩宠和慈祥与我们一同施行圣事。\par

\Emph{又：}\par

主要祝福我们，并使我们堪当常在祂的圣洁祭台前侍立，自今至永远，及世之世。\par

\Emph{又：}\par

愿受颂扬的是天主，祂在我们参与并事奉祂纯洁的奥迹时祝福并圣化我们众人，自今至永远，及世之世。\par

\Emph{执事诵念总祷文。}\par

XXII 让我们在平安中向主祈祷。\par

\Emph{众：}\par

主，求祢垂怜。\par

\Emph{执事：}\par

求祢拯救我们，垂怜我们，怜悯并保守我们，天主，藉祢的恩宠。\par

为那来自高天的平安、天主的慈祥和我们灵魂的救恩，\par

让我们恳求主。\par

为普世的平安和天主所有神圣教会的合一，\par

让我们恳求主。\par

为那些在天主神圣教会中结出果实、辛勤工作的人；为那些纪念穷人、寡妇和孤儿、客旅和贫困者的人；以及为那些请求我们在祈祷中纪念他们的人，\par

让我们恳求主。\par

为那些年老体衰、患病受苦、被不洁之灵困扰的人，为他们在天主那里的速效痊愈和救恩，\par

让我们恳求主。\par

为那些度守贞、独身和克修生活的人，以及那些在神圣婚姻中的人；并为我们那些在山岭、洞穴和地穴中战斗的圣父和弟兄们，\par

让我们恳求主。\par

为基督徒中航行、旅行、寄居的人，以及我们那些被俘、流放、监禁和苦役中的弟兄们，使他们平安归来，\par

让我们恳求主。\par

为赦免我们的罪过，宽恕我们的过犯，并使我们脱离一切患难、忿怒、危险、困迫和敌人的攻击，\par

让我们恳求主。\par

为天气适宜、及时降雨、恩泽甘露、五谷丰登、好季节的圆满结束，以及一年的冠冕，\par

让我们恳求主。\par

为我们现在和时时在这神圣时辰与我们一同祈祷的父老弟兄们——他们的热忱、辛劳和恳切，\par

让我们恳求主。\par

为每一个在患难和困苦中、需要天主慈悲和扶助的基督徒灵魂；为迷途者的归正、病者的痊愈、被俘者的释放、已亡父老弟兄的安息，\par

让我们恳求主。\par

为我们的祈祷在天主前蒙垂听和悦纳，并为祂丰富的仁慈和怜悯倾降于我们，\par

让我们恳求主。\par

并为这所呈献的、宝贵的、天上的、不可言喻的、纯洁的、光荣的、可畏的、可怕的、神圣的恩赐，以及那侍立并奉献它们的司铎的得救，\par

让我们向上主天主呈上恳求。\par

\Emph{众：}\par

主，求祢垂怜。\par

\Emph{（三次）}\par

\Emph{然后司铎在恩赐上画十字圣号，并站立，分别念诵如下：}\par

XXIII 光荣在天上归于天主，平安在地上，善意在人中间，等等。\par

\Emph{（三次）}\par

主，求祢开启我的口，我的口要颂扬祢的赞美。\par

\Emph{（三次）}\par

主，愿我的口充满祢的赞美，使我终日宣扬祢的光荣和威严。\par

\Emph{（三次）}\par

归于圣父。阿们。归于圣子。阿们。归于圣神。阿们。自今至永远，及世之世。阿们。\par

\Emph{然后他向左右俯首，说：}\par

XXIV. 你们要与我一同尊主为大，让我们一同高举祂的圣名。\par

\Emph{他们俯首回答说：}\par

圣神要降临到你身上，至高者的能力要荫庇你。\par

\Emph{然后司铎放长声音：}\par

主宰的主，祢曾以怜悯和慈悲眷顾我们，并赐予祢卑微、有罪、不配的仆人胆量，得以站立在祢的圣祭台前，并为我们的罪和百姓的过犯向祢献上这可畏的无血祭献。求祢垂顾祢无用的仆人，并因祢的慈悲涂抹我的过犯；并洁净我的嘴唇和心灵免受一切肉体和灵魂的玷污；并从我身上除去一切可耻和愚昧的思念，以祢至圣圣神的能力使我适于这项服务；并因祢的良善恩准我亲近祢的祭台。\par

主，也求祢悦纳我们手中所献的这些礼品，屈尊俯就我的软弱；不要将我逐离祢的圣容，也不要厌弃我的不配；但求祢按祢的大慈大悲怜悯我，按祢丰富的怜悯赦免我的过犯，好使我能无可指摘地来到祢的荣耀前，堪受祢独生子的保护，并蒙祢至圣圣神的光照，使我不致像罪奴般被驱逐，而能如同祢的仆人在祢面前获得恩宠、怜悯和罪赦，在今世和来世都是如此。\par

我恳求祢，全能的上主，大能的主，俯听我的祈祷；因为祢是在一切内运行一切的那一位，我们众人凡事都寻求来自祢和祢独生子，以及良善、赋予生命、同一体的圣神的帮助与支援，从现在直到永远。\par

二十五。天主，祢以伟大而不可言喻的爱，派遣祢的独生子来到世界，为要寻回迷失的羊，求祢不要转面不顾我们这些罪人，因为我们藉这可畏的无血祭献抓住祢；因为我们不倚靠自己的义德，而倚靠祢的良善仁慈，祢以此救赎了我们的族类。\par

我们恳求并哀求祢的良善，愿这救恩的奥迹由我们施行，为祢的子民不成为定罪，而是为得赦罪、为灵魂和身体的更新、为取悦祢——天主和圣父——因祢独生子的怜悯和慈爱，祢与祂同受赞颂，偕同祢至圣、良善、赋予生命的圣神，从现在直到永远。\par

二十六。主天主，祢创造了我们，并赐予我们生命，祢向我们显示了救恩的道路，祢赐予我们天上奥秘的启示，并藉祢至圣圣神的能力指派我们从事这项职务，求祢恩准，主上，使我们成为祢新约的仆人，祢纯洁奥迹的执事，并照着祢的大慈大悲接纳我们亲近祢的圣祭台，使我们堪当为我们的过犯和百姓的过犯向祢献上礼品和祭献；并求祢恩赐我们，主，以全然的敬畏和纯洁的良心向祢献上这属灵的无血祭献，并仁慈地接纳它，归入祢在天上的圣洁属灵祭台，作为馨香的属灵气味，求祢应允将祢至圣圣神的恩宠倾注在我们身上。\par

天主，求祢垂顾我们，并眷顾我们这合理的敬礼，接纳它，如同祢接纳了亚伯尔的礼品、诺厄的祭献、梅瑟和亚郎的司祭职、撒慕尔的和平祭、达味的忏悔、匝加利亚的馨香。正如祢从祢的宗徒手中接纳了这项真实的敬礼，求祢也照祢的良善从我们罪人手中接纳这些献上的礼品；并恩准我们的奉献蒙祢悦纳，由圣神所圣化，作为我们过犯和百姓错误的赎罪祭；也为那些先前安眠的灵魂；使我们这些祢卑微、有罪、不配的仆人，得以毫无诡诈地堪当侍奉祢的圣祭台，领受忠心而有智慧管家的赏报，并在祢公义良善报应的可畏之日寻获恩宠和怜悯。\par

\Emph{帷幔祈祷。}\par

二十七。我们感谢祢，主，我们的天主，因为祢赐予我们胆量进入祢的圣所，这圣所祢已藉祢基督血肉的帷幔为我们更新为一条又新又活的道路。因此，我们既堪当进入祢荣耀会幕的处所，进入帷幔之内，并瞻仰至圣所，就俯伏在祢的良善面前：\par

主，求祢垂怜我们，因为我们充满畏惧和战栗，正要站立在祢的圣祭台前，并为我们的罪和百姓的过犯献上这可畏的无血祭献。天主，求祢发出祢的良善恩宠，圣化我们的灵魂、身体和精神；并将我们的思念转向圣洁，使我们能以纯洁的良心向祢献上和平祭，即赞颂的祭献：\par

\Emph{（高声。）}\par

藉着祢独生子的怜悯和慈爱，祢与祂同受赞颂，偕同祢至圣、良善、赋予生命的圣神，从现在直到永远。\par

\Emph{民众。}\par

阿们。\par

\Emph{司铎。}\par

愿平安归于众人。\par

\Emph{执事。}\par

让我们恭敬地站立，让我们在敬畏天主和痛悔中站立；让我们专心于圣体共融礼仪，向天主献上平安祭。\par

\Emph{民众。}\par

平安的奉献，赞颂的祭献。\par

\Emph{司铎 [现在从饼酒的奉献物上揭去帷幔。]}\par

并向我们清楚地启示它，揭开那些以象征方式暗晦地笼罩着这神圣礼仪的帷幔。以绝对的光明充满我们理智的视觉，并洁净我们卑微之身脱离一切血肉和灵魂的玷污，使它堪当这可畏可惊的亲近：因为祢是至仁慈、全恩宠的天主，我们向祢献上赞颂和感谢，圣父、圣子、圣神，从现在直到永远。\par

% structure: p0219 source=h2 target=section reason=relative-heading-rank; base=h2

\section{3}

\Emph{感恩经。}\par

\Emph{然后他高声说：——}\par

二十八。主和圣父的慈爱，主和圣子的恩宠，圣神的共融与恩赐，与你们众人同在。\par

\Emph{民众。}\par

也与你的心灵同在。\par

\Emph{司铎。}\par

请举起你们的心意和心灵。\par

\Emph{民众。}\par

这是合理且理所当然的。\par

\Emph{然后司铎祈祷。}\par

确实，赞美祢、歌颂祢、称颂祢、朝拜祢、光荣祢、感谢祢，是合宜而正当的，适当而应当的；祢是万有——有形与无形受造物的创造者，永恒美善的宝藏，生命与不朽的泉源，万有的天主和主：\par

诸天诸天赞美祂，以及其中的所有天军；太阳、月亮和众星宿的合唱；大地、海洋和其中的一切；天上的耶路撒冷、天上的集会、以及写在诸天上的首生者的教会；义人和先知们的灵魂；殉道者和宗徒们的灵魂；天使、总领天使、上座者、宰制者、率领者、掌权者、可畏的德能；以及众目革鲁宾和六翼塞拉芬，它们用两个翅膀遮面，两个翅膀遮脚，两个翅膀飞翔，不停地以口唇相互呼喊，以不息的赞美：\par

\Emph{（高声朗读。）}\par

以高声歌唱祢威严光荣的凯旋赞美诗，高声呼喊、赞美、欢呼并说：—\par

\Emph{民众。}\par

\Quote{圣、圣、圣，万军的上主，天地充满祢的荣耀。欢呼之声于至高之天；奉主名而来的当受赞美。欢呼之声于至高之天。}\par

\Emph{司铎在祭品上画十字圣号，说：—}\par

二十九．祢是圣的，永生的君王，圣德的赐予者和主；祢的独生子、我们的主耶稣基督也是圣的，祢藉着祂创造了万有；祢的圣神也是圣的，祂探寻万事，连祢深奥的事也探寻，天主：祢是圣的，全能、全权、美善、可畏、仁慈、极其怜悯祢的受造物；祢用土造了人，按照祢的肖像和模样；祢赐予他乐园的喜乐；当他违背了祢的诫命、堕落时，祢没有忽视他，也没有遗弃他，善者啊，却如同慈父一般责罚他，藉着法律召唤他，藉着先知教导他；之后祢又派遣了祢的独生子、我们的主耶稣基督亲自来到世上，为要藉着祂的来临更新并恢复祢的肖像；\par

祂从天降下，由圣神和童贞天主之母玛利亚取得肉躯，并居住在我们中间，完成了我们种族救恩的经纶；祂即将在十字架上忍受祂自愿且赋予生命的死亡——祂是无罪的，为我们罪人，在祂被出卖的那一夜，不，更确切地说，祂为世界的生命和救恩而自我交出，\par

\Emph{然后司铎手捧面饼，说：—}\par

三十．祂用祂圣洁、纯洁、无玷、永不朽坏的手拿起面饼，举目向天，将饼呈献给祢——祂的天主和圣父，祝谢了，祝圣了，擘开了，并赐给了我们——祂的门徒和宗徒，说：—\par

\Emph{执事说：}\par

\Quote{为罪过的赦免和永生。}\par

\Emph{然后他高声说：—}\par

\Quote{你们拿去吃：这是我的身体，为你们而擘开，并赐下，为赦免罪过。}\par

\Emph{民众。}\par

阿们。\par

\Emph{然后他拿起杯，说：—}\par

同样，晚餐后，祂拿起杯，将酒和水搀和，举目向天，将杯呈献给祢——祂的天主和圣父，祝谢了，祝圣并降福了它，以圣神充满它，并将它赐给祂的门徒，说：\Quote{你们都喝这个；这是我的新约之血，为你们和众人倾流，并分施，为赦免罪过。}\par

\Emph{民众。}\par

阿们。\par

\Emph{司铎。}\par

\Quote{你们这样做，为纪念我；因为每当你们吃这个饼，喝这个杯，就是宣告主的死亡，并宣认祂的复活，直到祂来。}\par

\Emph{执事说：—}\par

\Quote{我们信并宣认：}\par

\Emph{民众。}\par

\Quote{我们宣告祢的死亡，主，并宣认祢的复活。}\par

\Emph{司铎（奉献礼）。}\par

三十一．因此，我们追念祂赋予生命的苦难、祂救赎的十字架、祂的死亡与埋葬、第三日从死者中复活、升天、坐在祢——我们的天主和圣父右边，以及祂第二次光荣而可畏的显现，那时祂必带着光荣来临，审判生者和死者，按照各人的行为施行报应；因此，我们这些罪人，向祢、主，奉献这可敬而不流血的祭献，恳求祢不要按我们的罪恶对待我们，也不要按我们的过犯报应我们；\par

但愿祢，按照祢的仁慈和不可言喻的恩爱，涂抹并擦去控告我们这些恳求者的罪状，而将祢属天永恒的恩赐（就是眼所未见、耳所未闻、人心所未想到的）赐给我们，就是祢为爱祢的人所预备的，天主；仁慈的主，不要因我或因我的罪而弃绝祢的子民：\par

\Emph{然后他说三次：—}\par

因为祢的子民和祢的教会向祢恳求。\par

\Emph{民众。}\par

\Quote{主、我们的天主、全能圣父，求祢垂怜我们。}\par

\Emph{司铎又说（呼求圣神祷文）：—}\par

三十二．恳求祢垂怜我们，全能的天主。\par

恳求祢垂怜我们，我们的救主天主。\par

恳求祢垂怜我们，天主，按照祢伟大的仁慈，并派遣祢的至圣圣神降临我们和这些奉献的礼品上。\par

\Emph{然后他低头说：—}\par

至高无上、赋予生命的圣神，祂与祢——我们的天主和圣父，并与祢的独生子同坐宝座，与祢一同为王；祂是同体、同永的；祂曾在法律、先知和祢的新约中发言；祂以鸽子的形状降临在约旦河我主耶稣基督身上，并停留在祂内；祂在五旬节日，在圣而光荣的熙雍楼房，以火舌的形状降临在祢的宗徒身上：祢的这位至圣圣神，求主，降临在我们和这些奉献的圣祭之上；\par

\Emph{他起身高声说：—}\par

愿祂藉着祂圣洁、美善而光荣的显现来临，圣化这饼，使之成为祢基督的圣体。\par

\Emph{民众。}\par

阿们。\par

\Emph{司铎。}\par

而这杯，成为祢基督的宝血。\par

\Emph{民众。}\par

阿们。\par

\Emph{司铎独自站立。}\par

三十三．使它们为所有领受者成为罪过的赦免和永生，成为灵魂和身体的圣化，成为善工果实的结出，成为祢圣而公教会的建立，这教会祢建立在信德之磐石上，使地狱的门不能胜过她；拯救她脱离一切异端和绊脚石，以及作恶的人，保守她直到时期的圆满。\par

\Emph{俯首后说：——}\par

三十四．我们也把这些祭品呈献给祢，主啊，为了圣所，就是祢藉着祢基督的神圣显现和祢至圣圣神的降临所荣耀的地方；特别是为了光荣的熙雍，众教会之母；并为全世界祢的圣而公宗徒的教会：就是现在，主啊，求祢将祢至圣圣神的丰厚恩赐赐予她。\par

主啊，也求祢纪念其中的圣教父和弟兄们，以及全世界的主教们，他们正确地分施祢真理的圣言。\par

主啊，也求祢纪念各城各地，以及住在其中的真信德之人，他们的平安与安全。\par

主啊，求祢纪念航海的、旅行的、寄居异乡的基督徒；我们的教父和弟兄们，他们在锁链、监禁、被掳和流放中；在矿坑中、受酷刑和苦役中。主啊，求祢纪念患病和受苦的，以及被不洁之灵困扰的，求祢，天主，赐他们速愈和救恩。\par

主啊，求祢纪念每个在困苦和艰难中的基督徒灵魂，需要祢的仁慈和援助，天主；以及迷途者的归正。\par

主啊，求祢纪念我们的教父和弟兄们，他们辛劳工作，为了祢的圣名而服事我们。\par

主啊，求祢以善纪念众人：主啊，求祢怜悯众人，求祢与众人和好：求祢赐祢百姓的群众平安：除去绊脚石：止息战争：使异端之兴起止息：求祢将祢的平安和祢的爱赐给我们，天主，我们的救主，地极之人的盼望。\par

主啊，求祢纪念风调雨顺、平安的甘霖、有益的露水、果实的丰盛，并以祢的良善冠冕年岁；因为众人都仰望祢，祢按时赐给他们食物：祢张开手，使一切有生命的都饱足喜乐。\par

主啊，求祢纪念那些在祢教会的圣所中结果实并光荣劳作的人；以及那些不忘穷人、寡妇、孤儿、旅客和贫困者的人；和所有希望我们在祈祷中纪念他们的人。\par

此外，主啊，求祢乐意纪念那些今日将这些祭品带到祢圣祭台的人，以及每个人为何奉献或怀着何种心意，还有那些刚刚被向祢提及的人。\par

主啊，求祢按祢丰盛的慈爱和怜悯，也纪念我，祢卑微无用的仆人；以及环绕祢圣祭台的执事们，并恩赐他们无瑕的生活，保守他们的服务不受玷污，为他们赢得良好的品位，使我们能同所有从创世以来蒙祢悦纳的圣人们——祖先、族长、先知、宗徒、殉道者、宣信者、教师、圣人，以及每个在祢基督的信德中得以成全的义人的灵魂——获得怜悯和恩宠，世世代代。\par

三十五．万福，玛利亚，满被圣宠者：主与尔同在；女中尔为赞美，尔胎子并为赞美，因为尔生了我等灵魂的救主。\par

\Emph{执事。}\par

三十六．主天主，求祢纪念我们。\par

\Emph{司铎俯首说：——}\par

主天主，求祢纪念我们曾提及和未曾提及的众灵魂和一切血肉，他们是真信德之人，从义人亚伯尔直到今日：求祢使他们在活人之地，在祢的国度，在乐园的喜乐中，在亚伯拉罕、依撒格、雅各伯——我们圣教父的怀抱中安息；那里没有痛苦、忧伤和哀叹：祢圣容的光辉照耀他们，并永远光照他们。\par

主啊，求祢使我们生命的终结成为基督徒的、可接纳的、无瑕的、平安的，主啊，求祢将我们聚集在祢选民脚下，随祢意愿，按祢心意；只求毫无羞愧和过犯，藉着祢的独生子，我们的主、天主和救主耶稣基督：因为惟有祂是没有罪而出现在地上的。\par

\Emph{执事。}\par

让我们祈祷：——\par

为全世界的和平与安定，以及天主的各圣教会，并为每个人奉献的目的或他所怀的意愿：并为周围站立的人群，以及所有的男人和女人：\par

\Emph{民众。}\par

并为所有的男人和女人。（阿们。）\par

\Emph{司铎高声说：——}\par

因此，为他们也为我们，求祢在祢的良善和爱中：\par

\Emph{民众。}\par

宽恕、赦免、饶恕，天主，我们的过犯，自愿的和非自愿的：行为上的和言语上的：明知的和不知的：黑夜的和白日的：思想中的和意图中的：在祢的良善和爱中，求祢宽恕我们这一切。\par

\Emph{司铎。}\par

藉着祢独生子的恩宠、怜悯和爱，祢与他，连同至圣、良善、赋予生命的圣神，同受赞颂和荣耀，自今而永远，及至无穷世。\par

\Emph{民众。}\par

阿们。\par

\Emph{司铎。}\par

三十七．平安与众人：\par

\Emph{民众。}\par

也与你的心灵。\par

\Emph{执事。}\par

再次，不断在平安中向主祈祷。\par

为献给主天主并已圣化的礼物，珍贵的、属天的、不可言喻的、纯洁的、荣耀的、可敬畏的、可怖的、神圣的；\par

让我们祈祷。\par

愿我们的主天主恩慈地将它们收纳到祂的祭台，那是圣洁的、高于诸天的、理性的、属灵的，作为馨香属灵的香气，为我们降下神圣的恩宠和至圣圣神的恩赐；\par

让我们祈祷。\par

已为信德的一致和祂至圣可敬圣神的共融祈祷之后；\par

让我们把自己和彼此，以及我们的整个生命，交托给基督我们的天主：\par

\Emph{民众。}\par

阿们。\par

\Emph{司铎祈祷。}\par

XXXVIII. 天主，我们的主、天主、救世主耶稣基督的圣父，荣耀的主，可赞美的本体，丰厚的良善，万有的天主和主宰，祢受赞颂直到永远，祢坐在革鲁宾之上，受色辣芬的光荣称颂；在祢面前有千千万万、万万千千的天使和总领天使的军队：祢收纳了奉到祢面前的礼品、奉献和初果，作为馨香的神怡之味，并乐于以祢的基督的恩宠和祢全圣神的临在，圣化它们、使之成全，良善之主啊，\par

主，求祢也圣化我们的灵魂、身体和心灵，触摸我们的理智，省察我们的良心，从我们身上驱逐一切邪恶的想象、一切不洁的感受、一切卑劣的欲望、一切不当的思想、一切嫉妒、虚荣和伪善、一切谎言、一切欺诈、一切世俗的情感、一切贪吝、一切虚浮的荣耀、一切冷漠、一切恶习、一切情欲、一切愤怒、一切怨恨、一切亵渎、一切不合祢圣意的肉身和心灵的冲动：\par

\Emph{(高声.)}\par

慈爱的主，求祢使我们堪当以坦然无惧的心、无罪的良心、纯洁的心灵、痛悔的精神、无愧的面容、圣化的嘴唇，敢于呼求祢——圣洁的天主、天上的圣父——并且说：\par

\Emph{众}\par

我们的天父，愿祢的名被尊为圣；等等直到光荣颂。\par

\Emph{司铎鞠躬，念（插祷）：——}\par

\Emph{不要让我们陷于诱惑，万军之主，祢知道我们的软弱，但救我们脱离那恶者及其作为，以及他的一切恶意和诡诈，因祢的圣名，这名已加在我们卑微的身上：}\par

\Emph{(高声.)}\par

因为国度、权能和光荣都是祢的，圣父、圣子、圣神，从今时直到永远。\par

\Emph{众}\par

阿们。\par

\Emph{司铎}\par

三十九。愿平安与众人。\par

\Emph{众}\par

也与你的心灵。\par

\Emph{执事}\par

让我们向主低头。\par

\Emph{众}\par

主，归于祢。\par

\Emph{司铎祈祷，如是说：——}\par

主，我们祢的仆人在祢的圣祭台前向祢低头，等候从祢而来的丰盛慈悲。\par

主，求祢向我们施放丰厚的恩宠和祝福；并圣化我们的灵魂、身体和心灵，使我们堪当成为祢神圣奥迹的领受者和分享者，以得罪赦和永生：\par

\Emph{(高声.)}\par

因为祢，我们的天主，以及祢的独生子，和祢的全圣神，是可钦崇和应受光荣的，从今时直到永远。\par

\Emph{众}\par

阿们。\par

\Emph{司铎高声说：——}\par

愿圣、同质、不受造、可钦崇的圣三的恩宠和慈悲，与我们一起。\par

\Emph{众}\par

也与你的心灵。\par

\Emph{执事}\par

怀着对天主的敬畏，让我们注意。\par

\Emph{司铎秘密地说：——}\par

圣洁的主，祢居于圣所，求祢以祢恩宠的圣言和祢全圣神的降临圣化我们；因为祢，主，曾说过：\Quote{你们应是圣的，因为我是圣的。}主，我们的天主，不可理解的天主圣言，与圣父和圣神同一实体、同永远、不可分，求祢在祢神圣而不流血的祭献中接纳这纯洁的赞歌；与革鲁宾和色辣芬，并从我——一个罪人——一同呼喊说：——\par

\Emph{他拿起祭品，高声说：——}\par

四十。圣物给圣人。\par

\Emph{众}\par

唯独一位是圣的，唯一的主耶稣基督，为光荣天主圣父，愿光荣归于祂，直到永远。\par

\Emph{执事}\par

\Emph{四十一。为赦免我们的罪过，并为我们的灵魂做赎罪，并为每一个在困苦和磨难中的灵魂，需要天主的慈悲和救助，并为迷途者的归正、病人的痊愈、俘虏的释放、已故先人和弟兄们的安息；}\par

让我们全体热切地说：主，垂怜：\par

\Emph{众（十二次）}\par

主，垂怜。\par

\Emph{然后司铎擘饼，将一半拿在右手，一半拿在左手，将右手所持的蘸入圣爵，说：——}\par

全圣身体和宝血——我们的主、天主、救世主耶稣基督——的结合。\par

\Emph{然后他在左手中的饼上画十字圣号；然后用那已画过的饼画另一半；然后立即开始擘开，并在一切之前，给每只圣爵各放一块，说：——}\par

它已成为一体，并被圣化、成全，因圣父、圣子、圣神之名，从今时直到永远。\par

\Emph{当他在饼上画十字圣号时，他说：——}\par

请看天主的羔羊，圣父之子，除免世罪者，为世界的生命和救恩而牺牲。\par

\Emph{当他给每只圣爵各放一块时，他说：——}\par

基督的一份圣物，充满恩宠和真理，来自圣父和圣神，愿光荣和能力归于祂，直到永远。\par

\Emph{然后他开始擘开，并说：——}\par

四十二。\BibleQuote{上主是我的牧者，我必一无所缺。}在青草地上，等等。\par

\Emph{然后}\par

\BibleQuote{我要时时称颂上主}，等等。\par

\Emph{然后}\par

\BibleQuote{我的天主，君王，我要尊崇祢}，等等。\par

\Emph{然后}\par

\BibleQuote{列邦，请赞美上主}，等等。\par

\Emph{执事}\par

主祭，请祝福。\par

\Emph{司铎}\par

上主将祝福我们，并使我们无可指摘，堪当领受祂纯洁恩赐的共融，从今时直到永远。\par

\Emph{当他们装满了，执事说：——}\par

主祭，请祝福。\par

\Emph{司铎说：——}\par

上主将祝福我们，并使我们堪能以纯洁的手指触摸那活炭，放在信友的口中，以净化并更新他们的灵魂和身体，从今时直到永远。\par

\Emph{然后}\par

\BibleQuote{请尝并看，上主是何等甘饴；}祂被分开而不分裂；分给信友而不减少；为赦免罪过，并获永生；从今时直到永远。\par

\Emph{执事}\par

在基督的平安内，让我们歌唱：\par

\Emph{歌咏者}\par

\BibleQuote{请尝并看，上主是何等甘饴。}\par

\Emph{司铎念领圣体前祷词。}\par

主，我们的天主，天上的食粮，宇宙的生命，我得罪了天，也得罪了祢，不配分享祢纯净的奥迹；但祢是仁慈的天主，求祢因祢的恩宠，使我能无罪地分享祢的圣体和宝血，以获得罪的赦免和永生。\par

\Emph{四十三、然后他将圣体分给圣职人员；当执事们拿起圣盘和圣爵准备分送给民众时，那拿着第一个圣盘的执事说：——}\par

主祭，请降福。\par

\Emph{司铎回答说：——}\par

光荣归于天主，祂已圣化我们且正在圣化我们众人。\par

\Emph{执事说：——}\par

天主，愿祢受举扬，高过诸天，愿祢的荣耀充满大地，愿祢的王国存留直到永远。\par

\Emph{当执事即将把它放在预备桌上时，司铎说：——}\par

愿主我们天主的圣名永受赞美。\par

\Emph{执事。}\par

怀着对天主的敬畏，并带着信德和爱德，上前来。\par

\Emph{民众。}\par

奉主名而来的，当受赞美。\par

\Emph{再者，当他将圣盘放在预备桌上时，他说：——}\par

主祭，请降福。\par

\Emph{司铎。}\par

天主，求祢拯救祢的子民，降福祢的产业。\par

\Emph{司铎再次说。}\par

光荣归于我们的天主，祂已圣化我们众人。\par

\Emph{当他将圣爵放回祭台上时，司铎说：——}\par

愿主的圣名永受赞美，直到永远。\par

\Emph{执事们和民众说：——}\par

主，求祢以赞美充满我们的口，以喜乐充满我们的唇，使我们终日歌颂祢的荣耀和伟大。\par

\Emph{又说：——}\par

我们感谢祢，基督我们的天主，因为祢使我们堪当分享祢的圣体和圣血，以获得罪的赦免和永生。我们恳求祢，因祢的良善和慈爱，使我们无罪地保存。\par

\Emph{最后一次进堂时的献香经。}\par

四十四、我们感谢祢，万有的救主和天主，感谢祢赐给我们的一切恩惠，并使我们分享祢神圣而纯洁的奥迹；我们向祢献上这乳香，恳求：将我们保护在祢翅膀的荫下，使我们直到最后一息仍堪当分享祢神圣的礼仪，以圣化我们的灵魂和身体，以继承天国产业；因为祢，天主，是我们的圣化者，我们向祢、圣父、圣子及圣神，献上赞美和感谢。\par

\Emph{执事在进堂时开始说。}\par

光荣归于祢，光荣归于祢，光荣归于祢，基督君王，圣父的独生圣言，因为祢使我们，祢有罪的、不配的仆人，堪当享受祢纯洁的奥迹，以获得罪的赦免和永生：光荣归于祢。\par

\Emph{当他完成了进堂，执事开始这样说：——}\par

四十五、再次、再次，并在一切时刻，在平安中，让我们恳求主。\par

愿我们对祂神圣礼仪的参与，能成为我们远离一切邪恶的助力，成为我们通往永生旅途的支撑，成为圣神的共融和恩赐；\par

让我们祈祷。\par

\Emph{司铎祈祷。}\par

纪念我们的至圣、纯洁、最荣耀、可赞美的夫人，天主之母和卒世童贞玛利亚，以及自创世以来所有蒙祢悦纳的诸圣，让我们将自己、彼此和我们的整个生命，交托给基督我们的天主。\par

\Emph{民众。}\par

主，归于祢。\par

\Emph{司铎。}\par

四十六、天主，祢以伟大而不可言喻的爱情，俯就祢仆人们的软弱，并使我们堪当分享这天上的宴席，不要因我们分享祢纯洁的奥迹而定我们有罪；但善良的主，求祢使我们常在圣神的圣化中，使我们成为圣洁，能在祢面容的光中，藉祢独生子、我们的主、天主和救主耶稣基督的仁慈，与自创世以来所有蒙祢悦纳的诸圣同分产业；祢与祢的至圣、善良、赋予生命的圣神同受赞颂：因为祢至宝贵而荣耀的圣名——圣父、圣子、圣神——现在、永远、直到无穷世，当受赞美和光荣。\par

\Emph{民众。}\par

阿们。\par

\Emph{司铎。}\par

愿平安与众人同在。\par

\Emph{民众。}\par

也与祢的心灵同在。\par

\Emph{执事。}\par

四十七、让我们向主低头。\par

\Emph{司铎。}\par

伟大而奇妙的天主，求祢垂顾祢的仆人，因为我们向祢低下了头。求祢伸出祢强而有力、满溢祝福的手，降福祢的子民。求祢保守祢的产业，使我们时时刻刻都能光荣祢，我们唯一永生的真天主，神圣而同一本体的圣三——圣父、圣子、圣神，现在、永远、直到无穷世。\par

\Emph{（高声。）}\par

因为赞美、尊荣、朝拜和感谢，都归于祢——圣父、圣子、圣神，现在、永远。\par

\Emph{执事。}\par

四十八、在基督的平安中，让我们歌唱：\par

\Emph{他又说：——}\par

在基督的平安中，让我们前行：\par

\Emph{民众。}\par

因主的名。主祭，请降福。\par

\Emph{执事念的遣散经。}\par

从光荣到光荣，我们赞美祢，我们灵魂的救主。光荣归于父、子及圣神，现在、永远、直到无穷世。我们赞美祢，我们灵魂的救主。\par

\Emph{司铎从祭台走向圣器室时念一篇祈祷文。}\par

四十九、我们力上加力，并在祢的圣殿中完成了全部神圣礼仪，即使现在，我们恳求祢，主、我们的天主，使我们堪当获得完美的慈爱；求祢修直我们的道路，将我们扎根在敬畏祢中，并使我们堪当天国，在基督耶稣我们的主内；祢与祢的至圣、善良、赋予生命的圣神同受赞颂，现在、始终、直到永远。\par

\Emph{执事。}\par

五十、再次、再次，并在一切时刻，在平安中，让我们恳求主。\par

\Emph{遣散后在圣器室念的祈祷文。}\par

主，祢在祢独生子、我们的主耶稣基督的至圣圣体和宝血的共融中，赐给了我们圣化；求祢也赐给我们祢善神的恩宠，使我们持守信德无瑕，引导我们达至完全的子女身份和救赎，以及将来的永恒喜乐；因为祢是我们的圣化和光明，天主，以及祢的独生子和祢的至圣圣神，现在、永远、直到无穷世。阿们。\par

\Emph{执事。}\par

在基督的平安中，让我们警醒。\par

\Emph{司铎}\par

\Quote{天主应受赞美，祂藉着圣洁、赋予生命、纯洁的奥迹的共融，祝福并圣化我们，自今至永远，世世无尽。阿们。}\par

\Emph{然后，赎罪祈祷。}\par

主耶稣基督，永生天主之子，羔羊与牧者，祢除免世罪，祢曾白白宽免两个债务人的债，并赐罪妇赦罪，又赐瘫子痊愈并赦免其罪；求祢宽恕、赦免、饶恕，天主，我们的过犯，无论有意无意，明知无知，因违命和不顺服所犯的，是祢至圣的圣神比祢仆人更知晓的：\par

如果属血肉、居住在这世界的人们，在任何事上偏离了祢的诫命，或受魔鬼怂恿，无论在言语或行为上，或如果他们遭受了诅咒，或由于某种特别的许愿，我恳求并祈求祢说不尽的慈爱，使他们能从自己的话中获得释放，从誓言和特别许愿中获得解脱，按祢的美善。\par

主，求祢俯听我为祢仆人的祈求，忽略他们的一切错误，不再记念；赦免他们每样过犯，无论有意无意，拯救他们脱离永罚：因为祢曾命令我们说：\BibleQuote{凡你们在地上所束缚的，在天上也要被束缚；凡你们在地上所释放的，在天上也要被释放。}因为祢是我们的天主，是能够怜悯、拯救和赦罪的天主；光荣归于祢，与永生之父及赋予生命的圣神，自今至永远，世世无尽。阿们。\par

\SourceRecord{Translated by James Donaldson.；Ante-Nicene Fathers；Vol. 7.；Edited by Alexander Roberts, James Donaldson, and A. Cleveland Coxe.；Buffalo, NY: Christian Literature Publishing Co.,；1886.；<http://www.newadvent.org/fathers/0717.htm>.}
